\section{Discussion and Next Writing Steps}

\subsection{Expected interpretation structure}
\begin{itemize}
  \item Start by separating the effects of replay, growth, and their combination; otherwise the replication may over-attribute improvements to neurogenesis when replay is responsible.
  \item Discuss performance and architecture size together because a dynamic-capacity method can improve accuracy by spending more parameters.
  \item Treat SD-19 results as the stronger external-validity check and MNIST results as the controlled replication anchor.
  \item Use ablations to identify whether the implementation conclusion depends on thresholding, growth geometry, global coupling, stopping rules, or training schedule.
\end{itemize}

\subsection{Limitations to state explicitly}
\begin{itemize}
  \item If only MNIST results are available, the replication should be described as partial and small-scale.
  \item If result variance across seeds is large, conclusions should be stated in terms of trends rather than definitive ranking.
  \item If NDL improves accuracy but grows substantially, the report should state the cost rather than only the gain.
  \item If the local implementation differs from the original paper in training schedule, replay implementation, or dataset preprocessing, those differences should be listed before interpreting discrepancies.
\end{itemize}

\subsection{Conversion from bullets to prose}
\begin{itemize}
  \item Convert each bullet into one short paragraph.
  \item Keep citations attached to the claim they support.
  \item Replace every result TODO marker with a table value, figure reference, or explicit statement that the experiment was not run.
  \item Replace any placeholder citation only after the corresponding Zotero entry has been verified in \artifact{references.bib}.
\end{itemize}
